% !TEX root = ../main.tex
\chapter{Allocator and Memory Accounting}
\label{ch:allocator}

Every tensor allocation in \Lucid\ passes through a custom size-class
allocator before reaching the underlying provider (\MLX's allocator
for GPU tensors, the system allocator for CPU tensors). The
allocator's job is twofold: amortise the cost of frequent
small-to-medium allocations by recycling buffers, and account for
memory in a way the user can query. This chapter describes both
roles.

The allocator is conventional in shape --- it is a size-class pool
of the variety familiar from CUDA-side ML frameworks --- but its
parameters are tuned to Apple silicon's memory characteristics and
its accounting separates pool state from provider state in a way
that has occasionally surprised users (\code{lucid.memory\_stats()}
can disagree with \code{Activity Monitor}; the explanation is in
\secref{sec:allocator:accounting}).

The chapter is organised bottom up: \secref{sec:allocator:why}
motivates the pool, \secref{sec:allocator:size_classes} describes
the size-class structure, \secref{sec:allocator:fast_path} describes
the allocation and free fast paths,
\secref{sec:allocator:accounting} describes the accounting API,
\secref{sec:allocator:empty_cache} describes the drain operation,
\secref{sec:allocator:oom} describes OOM handling, and
\secref{sec:allocator:coordination} discusses the lack of
coordination with the provider's own pool.

\section{Why a Pool}
\label{sec:allocator:why}

The naive design is to call the provider's allocator
(\MLX\ or system \code{malloc}) on every tensor allocation and
release. This works correctly but is slow in two ways. First, each
provider call has nontrivial overhead --- \MLX's allocator
synchronises with the GPU's command queue to know which buffers are
in flight, which costs a microsecond at least and sometimes more.
Second, the provider's own pool may have its own size classes that
do not match \Lucid's typical allocation sizes, causing
fragmentation and wasted memory.

A pool above the provider solves both problems. The pool retains a
freelist of recently-freed buffers, so a request that matches a
freelist entry returns in constant time without touching the
provider. The pool also rounds allocation sizes to its own size
classes, which gives it deterministic memory behaviour independent
of the provider's internal heuristics.

The cost is the pool itself: a small amount of code, a small amount
of bookkeeping memory, and the rare case where the pool retains a
buffer that the provider would have preferred to free.
\secref{sec:allocator:empty_cache} describes the user-visible
mitigation for that case.

\section{Size Classes}
\label{sec:allocator:size_classes}

\Lucid's pool uses 23 size classes arranged in a roughly geometric
progression. The exact sizes are chosen to match typical ML
allocation patterns; \tabref{tab:allocator:size_classes} gives the
full list.

\begin{table}[t]
\centering
\small
\begin{tabular}{rrrrl}
\toprule
Idx & Size & Idx & Size & Notes \\
\midrule
 0 & 16\,B    & 12 & 64\,KB   & \\
 1 & 32\,B    & 13 & 128\,KB  & \\
 2 & 64\,B    & 14 & 256\,KB  & \\
 3 & 128\,B   & 15 & 512\,KB  & \\
 4 & 256\,B   & 16 & 1\,MB    & \\
 5 & 512\,B   & 17 & 2\,MB    & \\
 6 & 1\,KB    & 18 & 4\,MB    & \\
 7 & 2\,KB    & 19 & 8\,MB    & \\
 8 & 4\,KB    & 20 & 16\,MB   & \\
 9 & 8\,KB    & 21 & 32\,MB   & \\
10 & 16\,KB   & 22 & 64\,MB   & largest pooled \\
11 & 32\,KB   & ---  & $>$64\,MB & passthrough to provider \\
\bottomrule
\end{tabular}
\caption[Allocator size classes.]{The 23 size classes in \Lucid's
allocator pool. Each class has a per-device freelist of depth at
most \(k_{\max} = 32\). Requests over 64\,MB bypass the pool and
go directly to the provider, on the rationale that very large
buffers are rarely reused at the same size and pooling them would
waste memory.}
\label{tab:allocator:size_classes}
\end{table}

\begin{itemize}
  \item Small classes (16, 32, 64, 128, 256, 512, 1024 bytes) cover
    per-scalar workspaces, small index arrays, and small parameter
    buffers.
  \item Medium classes (2K, 4K, 8K, 16K, 32K, 64K, 128K, 256K) cover
    moderate activation buffers and most layer weights.
  \item Large classes (512K, 1M, 2M, 4M, 8M, 16M, 32M, 64M) cover
    large activations, weight matrices, and gradient buffers.
\end{itemize}

A request for $n$ bytes is rounded up to the smallest size class
$\geq n$. Requests larger than 64\,MB are passed through to the
provider without pooling; the largest buffers are unlikely to be
reused at the same size and pooling them would waste memory.

The 23-class structure is determined empirically. We chose it after
profiling a representative set of training workloads (ResNet-50,
ViT-B, GPT-2-base) and observing the distribution of allocation
sizes; the geometric progression captures the long tail without
producing internal fragmentation greater than a factor of two on
typical requests.

\subsection{Internal fragmentation}
\label{ssec:allocator:internal_frag}

The trade-off in size-class pools is that a request for, say, 1100
bytes is rounded up to the 2K class, wasting 948 bytes. The worst
case is a request just over a size-class boundary, which wastes
just under one full size-class worth of memory.

For \Lucid's typical workloads the average wasted fraction is
roughly 15--20\,\%, which is acceptable given the speed-up from
pooling. The pool is not the right tool for memory-constrained
inference workloads where each byte matters; for those, the user
can disable the pool by calling \code{lucid.set\_allocator\_pool(False)}
and pay the provider-allocator overhead per call.

\subsection{External fragmentation}
\label{ssec:allocator:external_frag}

External fragmentation --- the inability to satisfy a request even
when the total free memory exceeds the request size --- is largely
a non-issue under unified memory on Apple silicon. The kernel's VM
subsystem is happy to reorganise pages on demand; the allocator
sees a flat address space and does not need to perform compaction.
This is one of the simplifying assumptions Apple silicon's memory
model affords.

\section{Freelist and Depth Limits}
\label{sec:allocator:freelists}

Each size class has a freelist: a stack of recently-freed buffers
of that size. The freelist is bounded by the constant
\code{kMaxDepth = 32}. When a buffer is freed:

\begin{itemize}
  \item If the freelist has fewer than 32 entries, the buffer is
    pushed onto the stack.
  \item Otherwise, the buffer is returned to the provider.
\end{itemize}

The depth limit is what keeps the pool's memory footprint bounded.
Without it, a training loop that allocates and frees many tensors
per iteration would accumulate freelist entries indefinitely.

The choice of 32 is empirical. We measured the freelist depth
distribution across representative workloads and observed that
depths above 8 rarely benefit from re-allocation (the same buffer
is not re-requested often enough), but the cost of evicting at
depth 32 versus depth 8 is small in the workload we measured. The
slightly conservative cap gives some head-room for workloads we did
not profile.

\subsection{The freelist as a stack, not a queue}
\label{ssec:allocator:lifo}

The freelist is LIFO: the most recently freed buffer is the first
candidate for reuse. This is the right policy for cache behaviour:
the most recently used buffer is most likely to still be resident
in SLC, so handing it back to a re-allocator is a cache-warm
operation. FIFO would have the opposite (cache-cold) behaviour.

The implementation is a \code{std::vector} used as a stack; push
and pop are \code{back()} operations.

\section{The Allocation Fast Path}
\label{sec:allocator:fast_path}

The allocation path, in pseudocode:

\begin{lstlisting}[style=lucidcpp,language=C++]
void* Allocator::alloc(size_t nbytes, Device dev) {
  size_t cls = size_class_index(nbytes);   // O(log 23) = O(1)
  Freelist& fl = freelists_[dev][cls];

  if (!fl.empty()) {
    void* p = fl.pop_back();
    stats_[dev].allocated += class_size(cls);
    stats_[dev].alloc_calls++;
    return p;
  }

  void* p = provider_alloc(class_size(cls), dev);
  if (!p) throw LucidError("out of memory");
  stats_[dev].allocated += class_size(cls);
  stats_[dev].reserved  += class_size(cls);
  stats_[dev].alloc_calls++;
  return p;
}
\end{lstlisting}

The fast path (freelist hit) is a few memory accesses and a stack
pop, on the order of 50 nanoseconds. The slow path
(\code{provider\_alloc}) ranges from microseconds (system
\code{malloc}) to tens of microseconds (\MLX's command-queue-aware
allocator). On a training step that allocates several thousand
buffers, the freelist hit rate is typically above 95\,\% after the
first dozen iterations.

\subsection{Thread safety}
\label{ssec:allocator:thread_safety}

The freelists are protected by a per-device mutex. The lock is held
only during the freelist operation itself; the provider call (if
needed) is performed without the lock. This is acceptable because
provider allocators are themselves thread-safe.

For \Lucid's typical workload --- a single training thread
dispatching to the pool --- the lock is uncontended and adds a few
nanoseconds. For multi-threaded data loaders that allocate from
worker threads, the lock can be contended; we measured a worst-case
overhead of 10\,\% on a workload with 8 simultaneous worker
threads, which we judge acceptable.

\section{The Deallocation Path}
\label{sec:allocator:dealloc}

The deallocation path:

\begin{lstlisting}[style=lucidcpp,language=C++]
void Allocator::free(void* p, size_t nbytes, Device dev) {
  size_t cls = size_class_index(nbytes);
  Freelist& fl = freelists_[dev][cls];

  stats_[dev].allocated -= class_size(cls);
  stats_[dev].free_calls++;

  if (fl.size() < kMaxDepth) {
    fl.push_back(p);
  } else {
    provider_free(p, dev);
    stats_[dev].reserved -= class_size(cls);
  }
}
\end{lstlisting}

The fast path (push to freelist) is again on the order of 50
nanoseconds. The slow path (return to provider) is the bottleneck
only when the freelist is consistently saturated.

\section{Memory Accounting}
\label{sec:allocator:accounting}

The pool maintains per-device statistics that the user can query.
The public Python API:

\begin{lstlisting}[style=lucid,language=LucidPython]
stats = lucid.memory_stats('metal')
# {'allocated': 4_812_345_678,
#  'reserved':  5_368_709_120,
#  'peak_allocated': 5_201_022_400,
#  'peak_reserved':  5_368_709_120,
#  'alloc_calls': 87234,
#  'free_calls':  82910}

lucid.reset_peak('metal')      # zero the peak high-water marks
\end{lstlisting}

The fields:

\begin{itemize}
  \item \code{allocated}: total bytes currently in use by live
    tensors (sum across size classes).
  \item \code{reserved}: total bytes the pool has from the
    provider (sum of buffers in use plus buffers on freelists).
    Always $\geq$ \code{allocated}.
  \item \code{peak\_allocated}: high-water mark of \code{allocated}
    since the process began or since \code{reset\_peak()}.
  \item \code{peak\_reserved}: corresponding high-water mark of
    \code{reserved}.
  \item \code{alloc\_calls}, \code{free\_calls}: cumulative
    counters since process start.
\end{itemize}

\subsection{Pool vs.\ provider stats}
\label{ssec:allocator:pool_vs_provider}

A common confusion: \code{lucid.memory\_stats()['allocated']} is
the pool's view of allocated memory, not the system's view of
memory used by the process. The provider may hold additional memory
in its own caches (\MLX's own allocator pool has separate
bookkeeping; the system allocator's free regions are similar), and
the kernel may have additional buffers allocated for the Metal
command queue itself.

The user-visible reconciliation:

\begin{itemize}
  \item \code{allocated} from \Lucid\ corresponds to live tensors.
  \item \code{reserved} from \Lucid\ minus \code{allocated} is
    held by \Lucid's freelists.
  \item Activity Monitor's reported memory minus \Lucid's
    \code{reserved} is held by \MLX, the Metal driver, and the
    rest of the process.
\end{itemize}

\code{lucid.empty\_cache()} closes the gap between
\code{reserved} and \code{allocated} by draining the freelists; it
does not affect the gap between \code{reserved} and Activity
Monitor's number, which is the responsibility of \MLX's own pool.

\section{The empty\_cache API}
\label{sec:allocator:empty_cache}

\code{lucid.empty\_cache()} drains every freelist back to the
provider:

\begin{lstlisting}[style=lucid,language=LucidPython]
lucid.empty_cache()              # both devices
lucid.empty_cache('cpu')         # CPU only
lucid.empty_cache('metal')       # GPU only
\end{lstlisting}

The operation iterates each size class on the requested device,
popping every freelist entry and calling \code{provider\_free} on
it. After the call, \code{reserved == allocated} and the pool is
in its leanest state.

\subsection{When to call}
\label{ssec:allocator:when_empty_cache}

The recommended pattern is to call \code{empty\_cache} at
transition points:

\begin{itemize}
  \item Between training and validation phases in a training loop.
    The two phases often have different working-set sizes; draining
    in between lets the next phase build its own freelists
    appropriately.
  \item After loading a checkpoint that allocates and frees large
    temporary buffers. The pool may have retained buffers that the
    next training step will not need.
  \item Before sampling a memory snapshot for profiling, to ensure
    the snapshot reflects live state rather than freelist headroom.
\end{itemize}

\code{empty\_cache} is not free: each freelist drain is one
provider call, and a large pool can have several hundred entries
across size classes. The cost is on the order of milliseconds and
should not be called inside a tight loop.

\subsection{What empty\_cache does not do}
\label{ssec:allocator:not_empty_cache}

It does not affect live tensors. Anything the user holds a
reference to remains allocated.

It does not affect the autograd graph. Tensors retained by
\code{grad\_fn} chains for backward computation remain alive; the
graph must be released (by completing or releasing the forward
pass) before its retained tensors can be freed.

It does not free memory held by \MLX's own pool. \MLX\ has its own
\code{mx.clear\_memory()}-equivalent function; \Lucid\ does not
invoke it because the policies are independent. A user who wants
both layers drained should call both.

\section{Out-of-Memory Handling}
\label{sec:allocator:oom}

When the provider's allocator returns null, \Lucid's allocator
raises \code{LucidError("out of memory")} with the requested size,
the current pool state, and the device. The error message is
designed to be diagnostic:

\begin{verbatim}
LucidError: out of memory
  request: 268435456 bytes (256 MB) on device=metal
  pool reserved: 14782 MB
  pool allocated: 14503 MB
  pool peak reserved: 14782 MB
  hint: call lucid.empty_cache() to drain freelists;
        if still failing, reduce batch size or activation memory.
\end{verbatim}

The error is fatal in the sense that the framework does not
attempt a fallback: it does not silently reduce batch size, does
not silently move to CPU, does not silently retry with smaller
shapes. The reasoning is that any such fallback would change the
numerics of the training run in ways the user did not request, and
silent gradient changes are worse than an explicit failure.

\subsection{Memory pressure events versus OOM}
\label{ssec:allocator:pressure_oom}

OOM is the terminal failure mode. macOS memory-pressure events
(\code{DISPATCH\_MEMORYPRESSURE\_WARN},
\code{DISPATCH\_MEMORYPRESSURE\_CRITICAL}) are pre-OOM signals from
the kernel asking processes to release caches. \Lucid\ does not
itself subscribe to these events at the engine layer; the providers
(\MLX, \Accelerate) subscribe on their own behalf. The user can
respond by calling \code{empty\_cache} manually.

\section{Coordination with the Provider's Pool}
\label{sec:allocator:coordination}

\Lucid's pool sits above the provider's pool. The two pools have
independent freelists, independent size-class structures, and
independent eviction policies. There is no API by which \Lucid\
can ask \MLX\ to drop its own cache, nor vice versa.

The consequence is that the layered pools can together hold more
memory than either would alone. In the worst case, \Lucid's pool
holds a freelist entry that \MLX\ would have evicted, and
\MLX's pool holds the underlying buffer that \Lucid's freelist
just freed. We have not measured this composite overhead but
estimate it at 5--10\,\% of total reserved memory in typical
training workloads.

The simplification we adopted, for the sake of code clarity, was to
let each pool manage itself. Coordinating them would have required
a callback interface that we judged not worth the maintenance cost.

\subsection{When to disable Lucid's pool}
\label{ssec:allocator:disable_pool}

For memory-constrained workloads on small machines (8\,GB unified
memory or less), the layered overhead can matter. The user can
disable \Lucid's pool entirely:

\begin{lstlisting}[style=lucid,language=LucidPython]
lucid.set_allocator_pool(False)
# Now every alloc/free goes directly to the provider.
# The freelist depth is effectively zero.
\end{lstlisting}

The cost is the per-call overhead of provider allocation, which is
several microseconds. On a training step that allocates a few
thousand buffers, the disabled-pool overhead is around 10\,ms per
step --- noticeable but acceptable for workloads where memory is
tighter than time.

The default is enabled; disabling is the rare case.

\section{Summary and Forward References}
\label{sec:allocator:summary}

\Lucid's allocator is a size-class pool with 23 classes and a
per-class freelist of depth 32. It sits above the provider
(\MLX\ for GPU, the system allocator for CPU) and amortises
per-call overhead by recycling recently-freed buffers. The pool's
state is exposed through \code{lucid.memory\_stats} and drained
through \code{lucid.empty\_cache}; OOM raises \code{LucidError}
without silent fallback.

The next chapter, \chref{ch:backend_dispatch}, treats the
dispatcher that selects between the CPU and GPU backends.
\chref{ch:bridge_boundaries} catalogues the six places where the
allocator's contract crosses into external libraries.
\chref{ch:profiler_determinism} returns to the allocator from the
profiling angle, since per-op memory accounting depends on the
pool's bookkeeping.

\begin{lucidtip}
For the user, the two operational levers exposed by this chapter
are \code{lucid.memory\_stats()} (for observing) and
\code{lucid.empty\_cache()} (for acting). Both are cheap to call
during development and expensive enough not to call in a tight
loop. The OOM error message is designed to suggest \code{empty\_cache}
first, then a batch-size reduction.
\end{lucidtip}
